\documentclass[11pt,a4paper]{article}

% --- Πακέτα για Ελληνικά και Γραμματοσειρές μέσω XeLaTeX ---
\usepackage{fontspec}
\setmainfont{Times New Roman}
\usepackage{xgreek}

% --- Βασικά Πακέτα ---
\usepackage[margin=1in]{geometry}
\usepackage{amsmath, amssymb}
\usepackage{booktabs}
\usepackage{hyperref}
\usepackage{microtype}
\usepackage{float}

% --- Διαμόρφωση Εμφάνισης Κώδικα ---
\usepackage{listings}
\usepackage{xcolor}
\definecolor{codegreen}{rgb}{0,0.6,0}
\definecolor{codegray}{rgb}{0.5,0.5,0.5}
\definecolor{backcolour}{rgb}{0.95,0.95,0.92}

\lstset{
	backgroundcolor=\color{backcolour},   
	commentstyle=\color{codegreen},
	numberstyle=\tiny\color{codegray},
	basicstyle=\ttfamily\small,
	breakatwhitespace=false,         
	breaklines=true,                 
	captionpos=b,                    
	keepspaces=true,                 
	numbers=left,                    
	numbersep=5pt,                  
	showspaces=false,                
	showstringspaces=false,
	showtabs=false,                  
	tabsize=4
}

\title{\textbf{Υλοποιητική Άσκηση - Τεχνικές Εξόρυξης Δεδομένων}}
\author{\textbf{Χρήστος Μάριος Νικολόπουλος - 1100644}}
\date{Εαρινό Εξάμηνο 2025-2026}

\begin{document}
	
	\maketitle
	\newpage
	\tableofcontents
	\newpage
	
	% =========================================================================
	\section{Περιβάλλον Υλοποίησης \& Εγκατάσταση}
	
	\subsection{Γλώσσα και Αρχιτεκτονική Περιβάλλοντος}
	Σύμφωνα με τις προδιαγραφές της εργαστηριακής άσκησης, ως γλώσσα υλοποίησης ορίστηκε η \textbf{Python} (έκδοση $\ge$ 3.10). Για την ανάπτυξη και τη διασφάλιση της σταθερότητας του κώδικα, δημιουργήθηκε ένα απομονωμένο εικονικό περιβάλλον (\textit{virtual environment - .venv}). Αυτή η πρακτική αποτρέπει τις διενέξεις μεταξύ των εκδόσεων των βιβλιοθηκών του συστήματος και εξασφαλίζει την πλήρη αναπαραγωγιμότητα των πειραμάτων.
	
	\subsection{Αναλυτική Καταγραφή Βιβλιοθηκών Λογισμικού}
	Για την κάλυψη των αναγκών των τριών ερωτημάτων (EDA, Classification, Clustering), επιλέχθηκαν και χρησιμοποιήθηκαν οι εξής εξειδικευμένες βιβλιοθήκες:
	
	\begin{itemize}
		\item \textbf{\texttt{pandas} \& \texttt{numpy}:} Αποτελούν τον πυρήνα της διαχείρισης δεδομένων. Χρησιμοποιήθηκαν για τη φόρτωση των CSV αρχείων, τον εντοπισμό και την αφαίρεση διπλότυπων, τη διαχείριση ελλιπών τιμών, καθώς και για όλους τους απαραίτητους μετασχηματισμούς των πινάκων (DataFrames).
		\item \textbf{\texttt{scikit-learn} (\texttt{sklearn}):} Η κεντρική βιβλιοθήκη μηχανικής μάθησης. Μέσω αυτής υλοποιήθηκαν οι τεχνικές κανονικοποίησης, οι αλγόριθμοι επιβλεπόμενης μάθησης (\texttt{LogisticRegression}, \texttt{DecisionTreeClassifier}, \texttt{RandomForestClassifier}), η βελτιστοποίηση υπερπαραμέτρων (\texttt{GridSearchCV}, \texttt{StratifiedKFold}), καθώς και οι αλγόριθμοι ομαδοποίησης (\texttt{KMeans}, \texttt{AgglomerativeClustering}, \texttt{DBSCAN}).
		\item \textbf{\texttt{matplotlib} \& \texttt{seaborn}:} Βιβλιοθήκες οπτικοποίησης. Χρησιμοποιήθηκαν για τη σχεδίαση των πινάκων συσχέτισης (Heatmaps) και των πινάκων σύγχυσης (Confusion Matrices), οι οποίοι αποθηκεύονται αυτόματα σε μορφή εικόνας (\texttt{.png}).
	\end{itemize}
	
	\subsection{Βήματα Εγκατάστασης \& Εκτέλεσης}
	Οι εξαρτήσεις του έργου έχουν καταγραφεί συγκεντρωτικά στο αρχείο \texttt{requirements.txt}. Η διαδικασία στησίματος του περιβάλλοντος από το τερματικό περιλαμβάνει τα εξής βήματα:
	
	\begin{lstlisting}[language=bash, caption={Εντολές Εγκατάστασης και Εκτέλεσης Pipeline}]
		cd path/to/Data_mining_project/src
		
		python -m venv .venv
		
		.venv\Scripts\activate
		
		pip install --upgrade pip
		pip install -r requirements.txt
		
		python src/pipeline.py
	\end{lstlisting}
	
	\newpage
	% =========================================================================
	\section{Διαδικασία Υλοποίησης (Architecture \& Pipeline)}
	Η αρχιτεκτονική του συστήματος βασίζεται σε αντικειμενοστραφή σχεδιασμό (Object-Oriented Programming). Κάθε στάδιο του πειράματος έχει απομονωθεί σε ξεχωριστό αρχείο κώδικα (module), ενώ η συνολική ροή ενορχηστρώνεται από ένα κεντρικό pipeline αρχείο (\texttt{pipeline.py}).
	
	\subsection{Στάδιο 1: Φόρτωση Δεδομένων (\texttt{DataLoader})}
	Η αρχική ανάγνωση των δεδομένων υλοποιείται μέσω της κλάσης \texttt{DataLoader} στο αρχείο \texttt{data\_loader.py}. Η μέθοδος \texttt{load\_data()} εντοπίζει δυναμικά τον κατάλογο εργασίας και χρησιμοποιεί τη βιβλιοθήκη \texttt{glob} για να συλλέξει όλα τα αρχεία μορφής \texttt{.csv} από το μονοπάτι \texttt{../data/raw}. Για κάθε αρχείο, πραγματοποιείται αφαίρεση κενών διαστημάτων από τους τίτλους των στηλών (\texttt{columns.str.strip()}) για την αποφυγή σφαλμάτων αναφοράς, και όλα τα επιμέρους DataFrames ενοποιούνται με την \texttt{pd.concat()}.
	
	\begin{lstlisting}[language=Python, caption={Υλοποίηση Φόρτωσης Δεδομένων στην DataLoader}]
		def load_data(self):
		current_dir = os.getcwd()
		path = os.path.join(current_dir, "..", "data", "raw")
		all_files = glob.glob(os.path.join(path, "*.csv"))
		print("Loading Data")
		list_df = []
		for f in all_files:
		temp_df = pd.read_csv(f, low_memory=False)
		temp_df.columns = temp_df.columns.str.strip()
		list_df.append(temp_df)
		self.df = pd.concat(list_df, ignore_index=True)
		return self.df
	\end{lstlisting}
	
	\subsection{Στάδιο 2: Προεπεξεργασία \& Καθαρισμός (\texttt{Preprocess})}
	Η κλάση \texttt{Preprocess} (\texttt{preprocess.py}) αναλαμβάνει την εκκαθάριση του dataset πριν από τη μοντελοποίηση:
	\begin{itemize}
		\item \textbf{Διαχείριση Ελλιπών \& Άπειρων Τιμών:} Οι τιμές απείρου (\texttt{inf}, \texttt{-inf}) αντικαθίστανται με \texttt{NaN}. Οι στήλες \texttt{Flow Bytes/s} και \texttt{Flow Packets/s} συμπληρώνονται με τη μέση τιμή τους (\texttt{fillna}), ενώ τυχόν εναπομείναντα κενά διαγράφονται οριστικά (\texttt{dropna}).
		\item \textbf{Αφαίρεση Διπλοτύπων:} Η μέθοδος \texttt{removing\_dupl()} απομακρύνει τις πανομοιότυπες εγγραφές, θωρακίζοντας τα μοντέλα από το φαινόμενο του overfitting.
		\item \textbf{Επιλογή Χαρακτηριστικών (Feature Selection):} Υπολογίζεται ο πίνακας συσχέτισης (\texttt{corr().abs()}). Με κατώφλι \texttt{threshold=0.95}, η μέθοδος εντοπίζει και αφαιρεί τις στήλες του άνω τριγωνικού πίνακα που παρουσιάζουν υπερβολική πολυγραμμικότητα, μειώνοντας τη διαστασιμότητα.
	\end{itemize}
	
	\begin{lstlisting}[language=Python, caption={Μέθοδος Feature Selection βάσει Συσχέτισης}]
		def feature_selection(self, threshold=0.95):
		numeric_df = self.df.select_dtypes(include=np.number)
		corr_matrix = numeric_df.corr().abs()
		upper = corr_matrix.where(
		np.triu(np.ones(corr_matrix.shape), k=1).astype(bool)
		)
		to_drop = [col for col in upper.columns if any(upper[col] > threshold)]
		self.df.drop(columns=to_drop, inplace=True)
		print(f"\nFeature selection (threshold={threshold}): removed {len(to_drop)} columns")
		return to_drop
	\end{lstlisting}
	
	\subsection{Στάδιο 3: Επιβλεπόμενη Μάθηση (\texttt{SupervisedLearning})}
	Η κλάση \texttt{SupervisedLearning} (\texttt{modeling.py}) διαχειρίζεται την εκπαίδευση των διεργασιών ταξινόμησης.
	\begin{enumerate}
		\item \textbf{Διαχωρισμός \& Κανονικοποίηση:} Στο \texttt{pipeline.py} απομονώνεται στρωματοποιημένο δείγμα \textbf{50.000 εγγραφών}. Η \texttt{prepare\_data()} χωρίζει το δείγμα σε Train (60\%), Validation (20\%) και Test (20\%) με \texttt{stratify=y}, και ακολουθεί τυποποίηση των χαρακτηριστικών με \texttt{StandardScaler}.
		\item \textbf{Στρατηγική Grid Search:} Για τα μοντέλα Logistic Regression, Decision Tree και Random Forest (με \texttt{class\_weight='balanced'}), εκτελούνται δύο σενάρια: το Scenario 1 (χρήση \texttt{PredefinedSplit} για αξιολόγηση στο Validation set) και το Scenario 2 (κλασική διασταυρούμενη επικύρωση 5-fold CV στο Train set).
	\end{enumerate}
	
	\begin{lstlisting}[language=Python, caption={Υλοποίηση των δύο Σεναρίων Grid Search}]
		def run_classification(self, option):
		classifier = self.select_classifier(option)
		params     = self.get_grid(option)
		model_name = MODEL_NAMES[option]
		
		# Scenario 1 - Predefined split (Train/Val Split)
		split      = [-1] * len(self.X_train) + [0] * len(self.X_val)
		pds        = PredefinedSplit(test_fold=split)
		X_combined = np.concatenate((self.X_train, self.X_val), axis=0)
		y_combined = np.concatenate((self.y_train, self.y_val), axis=0)
		
		grid_split = GridSearchCV(classifier, params, cv=pds, scoring="f1_weighted", n_jobs=-1)
		grid_split.fit(X_combined, y_combined)
		best_split = self.evaluate(grid_split, "Scenario 1 - Train/Val Split", option)
		
		# Scenario 2 - 5 fold cross validation
		grid_cv = GridSearchCV(classifier, params, cv=5, scoring='f1_weighted', n_jobs=-1)
		grid_cv.fit(self.X_train, self.y_train)
		best_cv = self.evaluate(grid_cv, "Scenario 2 - 5-fold CV", option)
	\end{lstlisting}
	
	\subsection{Στάδιο 4: Μη Επιβλεπόμενη Μάθηση (\texttt{UnsupervisedLearning})}
	Η κλάση \texttt{UnsupervisedLearning} (\texttt{clustering.py}) εκτελείται σε υποσύνολο \textbf{5.000 δειγμάτων} χωρίς τη χρήση του \texttt{Label}. 
	Πραγματοποιείται εύρεση βέλτιστου $k$ (Elbow method \& Silhouette analysis) για τον K-Means, εκτέλεση Ιεραρχικής Ομαδοποίησης με διάφορα κριτήρια σύνδεσης (\texttt{ward}, \texttt{complete}, \texttt{average}), και Grid Search για τον DBSCAN με μεταβολή των παραμέτρων \texttt{eps} και \texttt{min\_samples}.
	
	\begin{lstlisting}[language=Python, caption={Διαδικασία DBSCAN Grid Search}]
		def dbscan_grid_search(self, eps_values=(0.3, 0.5, 1.0, 2.0), min_samples_values=(5, 10, 20)):
		results = []
		for eps in eps_values:
		for ms in min_samples_values:
		model  = DBSCAN(eps=eps, min_samples=ms, n_jobs=-1)
		labels = model.fit_predict(self.X)
		n_clusters = len(set(labels)) - (1 if -1 in labels else 0)
		
		if n_clusters >= 2:
		sil = silhouette_score(self.X, labels, sample_size=10000, random_state=42)
		else:
		sil = -1.0
		results.append({'eps': eps, 'min_samples': ms, 'silhouette': sil})
		best = max(results, key=lambda r: r['silhouette'])
		return self.run_dbscan(eps=best['eps'], min_samples=best['min_samples']), best
	\end{lstlisting}
	
	\newpage
	% =========================================================================
	% =========================================================================
	\section{Σχολιασμός Τελικών Αποτελεσμάτων (Evaluation \& Discussion)}
	
	\subsection{Προετοιμασία Δεδομένων \& Ανισορροπία Κλάσεων}
	Το αρχικό σύνολο δεδομένων (CIC-IDS-2017) χαρακτηρίζεται από τεράστιο όγκο ($2.830.743$ εγγραφές και $79$ χαρακτηριστικά). Η εφαρμογή του φιλτραρίσματος για την αφαίρεση άπειρων ($inf$) και ελλιπών ($NaN$) τιμών, σε συνδυασμό με τη μέθοδο επιλογής χαρακτηριστικών (Feature Selection) με κατώφλι συσχέτισης $\text{threshold} = 0.95$, οδήγησε στην αφαίρεση $23$ περιττών στηλών, αφήνοντας $55$ τελικά χαρακτηριστικά για τη μοντελοποίηση.
	
	Το πιο κρίσιμο στοιχείο που καθορίζει τη συμπεριφορά όλων των μοντέλων είναι η \textbf{ακραία ανισορροπία κλάσεων (Extreme Class Imbalance)}. Η κλάση \texttt{0} (BENIGN) κυριαρχεί καθολικά με πάνω από $2$ εκατομμύρια δείγματα, ενώ μειονοτικές κλάσεις όπως η \texttt{8} (Heartbleed) εμφανίζουν μόλις $11$ δείγματα σε ολόκληρο το dataset. Για τον λόγο αυτό, η στρωματοποιημένη δειγματοληψία $50.000$ γραμμών ήταν απαραίτητη για να καταστεί υπολογιστικά εφικτή η διαδικασία εκπαίδευσης.
	
	\subsection{Επιβλεπόμενη Μάθηση (Supervised Learning)}
	Η αξιολόγηση των τριών ταξινομητών πραγματοποιήθηκε σε test set $10.000$ δειγμάτων, συγκρίνοντας τη συμπεριφορά τους στο Scenario 1 (Train/Val Split) και στο Scenario 2 (5-fold Cross Validation).
	
	\subsubsection{Logistic Regression (Λογιστική Παλινδρόμηση)}
	Η Λογιστική Παλινδρόμηση παρουσίασε τη χαμηλότερη συνολική απόδοση με συνολική ακρίβεια $\text{Accuracy} = 0.89$ και $\text{Macro Avg } F_1\text{-Score} = 0.53$. Λόγω της γραμμικής της φύσης, αδυνατεί να διαχωρίσει ικανοποιητικά τις μειονοτικές κλάσεις:
	\begin{itemize}
		\item Στην κλάση \texttt{1}, εμφάνισε εξαιρετικά χαμηλό $\text{Precision} = 0.01$ αλλά υψηλό $\text{Recall} = 0.86$ (ή $0.71$ στο CV). Αυτό υποδηλώνει ότι το μοντέλο υπερ-προβλέπει τη συγκεκριμένη επίθεση, παράγοντας τεράστιο αριθμό ψευδώς θετικών αποτελεσμάτων (False Positives).
		\item Στις πολύ σπάνιες κλάσεις (\texttt{12}, \texttt{14}), το $F_1\text{-score}$ κυμάνθηκε κοντά στο μηδέν ($0.03$).
		\item Τα αποτελέσματα μεταξύ των δύο σεναρίων ήταν σχεδόν πανομοιότυπα με τις ίδιες βέλτιστες παραμέτρους ($\text{C}=10, \text{max\_iter}=1000$), αποδεικνύοντας ότι το μοντέλο έχει σταθερή συμπεριφορά αλλά πάσχει από \textit{underfitting}.
	\end{itemize}
	
	\subsubsection{Decision Tree \& Random Forest (Δενδρικά Μοντέλα)}
	Αντίθετα με τη γραμμική προσέγγιση, οι δενδρικοί αλγόριθμοι επέδειξαν \textbf{εξαιρετική συμπεριφορά} αγγίζοντας το $\text{Accuracy} = 1.00$ και $\text{Weighted Avg } F_1\text{-Score} = 1.00$.
	\begin{itemize}
		\item Το \textbf{Decision Tree} πέτυχε κορυφαίες επιδόσεις με τις μειονοτικές κλάσεις (π.χ. κλάση \texttt{1} με $F_1 = 0.80$ και κλάση \texttt{12} με $F_1 = 0.73$ στο Scenario 1). Η χρήση του κριτηρίου \texttt{entropy} με \texttt{max\_depth=20} προσέφερε την καλύτερη ισορροπία. Στο 5-fold CV, το μοντέλο χωρίς περιορισμό βάθους (\texttt{None}) διατήρησε την υψηλή του απόδοση ($\text{Macro } F_1 = 0.82$).
		\item Το \textbf{Random Forest} (με βέλτιστες παραμέτρους $200$ δέντρα και $\text{max\_depth}=20$ στο CV) επιβεβαίωσε την ανωτερότητά του ως ensemble μέθοδος. Εμφάνισε εξαιρετικό $F_1$-score $\ge 0.92$ για σχεδόν όλες τις κλάσεις επιθέσεων. 
		\item \textbf{Εξαίρεση στην κλάση 14:} Λόγω του ελάχιστου support ($3$ δείγματα), κανένας από τους αλγορίθμους δεν κατάφερε να την εντοπίσει σωστά ($F_1 = 0.00$), καθώς το μέγεθος του δείγματος δεν επαρκούσε για την εξαγωγή κανόνων.
	\end{itemize}
	\textbf{Συμπέρασμα Ταξινόμησης:} Τα δενδρικά μοντέλα είναι ιδανικά για την ανίχνευση δικτυακών εισβολών, καθώς οι κυβερνοεπιθέσεις συχνά αφήνουν διακριτά «αποτυπώματα» που εκφράζονται μέσω αυστηρών ιεραρχικών κατωφλίων (if-then rules) στα χαρακτηριστικά της ροής (π.χ. συγκεκριμένα Destination Ports ή Flow Durations).
	
	\subsection{Μη Επιβλεπόμενη Μάθηση (Unsupervised Learning \& Clustering)}
	Η συσταδοποίηση πραγματοποιήθηκε σε $5.000$ δείγματα χωρίς τη χρήση της μεταβλητής \texttt{Label}. Η σύγκριση των αποτελεσμάτων με τις πραγματικές ετικέτες αναδεικνύει τη γεωμετρική δομή του χώρου.
	
	\subsubsection{K-Means \& Ιεραρχική Ομαδοποίηση (Agglomerative)}
	Η ανάλυση Silhouette υπέδειξε ως βέλτιστο αριθμό συστάδων το $k=2$ με ένα εξαιρετικά υψηλό σκορ $\text{Silhouette} = 0.9752$ και πολύ χαμηλό δείκτη Davies-Bouldin ($\text{DB} = 0.0168$). 
	
	Εξετάζοντας τους πίνακες συσχέτισης των clusters με τα πραγματικά Labels (Confusion Matrices ομαδοποίησης), παρατηρούμε ότι τόσο ο \textbf{K-Means ($k=2$)} όσο και όλες οι παραλλαγές της \textbf{Ιεραρχικής Ομαδοποίησης (Ward, Complete, Average)} παρήγαγαν \textbf{ακριβώς το ίδιο αποτέλεσμα}:
	\begin{itemize}
		\item Το \textbf{Cluster 0} απορρόφησε σχεδόν το σύνολο των δειγμάτων ($4159$ BENIGN και σχεδόν όλες τις επιθέσεις).
		\item Το \textbf{Cluster 1} περιέχει \textbf{μόνο 1 δείγμα} (το οποίο ανήκει στην κλάση \texttt{0}).
	\end{itemize}
	Αυτή η συμπεριφορά, σε συνδυασμό με το γεγονός ότι η αύξηση των συστάδων σε $k=3$ έριξε κατακόρυφα το Silhouette score στο $0.4414$, υποδηλώνει την ύπαρξη ενός \textbf{ακραίου απομακρυσμένου σημείου (extreme outlier)} στα δεδομένα. Οι αλγόριθμοι που βασίζονται σε αποστάσεις (όπως ο K-Means και ο Hierarchical) αναγκάζονται να δημιουργήσουν ένα ξεχωριστό cluster μόνο για αυτό το outlier σημείο προκειμένου να ελαχιστοποιήσουν το σφάλμα, αποτυγχάνοντας να διαχωρίσουν τις πραγματικές επιθέσεις από τη φυσιολογική κίνηση. Επιπλέον, η PCA έδειξε ότι οι 2 πρώτες κύριες συνιστώσες ερμηνεύουν το $41.07\%$ της συνολικής διασποράς, επιβεβαιώνοντας ότι ο χώρος έχει υψηλή πολυπλοκότητα.
	
	\subsubsection{DBSCAN}
	Ο DBSCAN, ως αλγόριθμος βασισμένος στη πυκνότητα, συμπεριφέρθηκε τελείως διαφορετικά. Η βέλτιστη παραμετροποίηση που προέκυψε από το Grid Search ήταν $\text{eps}=2.0$ και $\text{min\_samples}=5$ ($\text{Silhouette} = 0.3109$).
	
	Ο αλγόριθμος εντόπισε $28$ διαφορετικές συστάδες και απομόνωσε $414$ σημεία ως θόρυβο (\texttt{Noise points}). Ο έλεγχος του πίνακα συσχέτισης δείχνει ότι ο DBSCAN κατάφερε να «σπάσει» τη μάζα των δεδομένων σε πολλές μικρές, πυκνές υποομάδες:
	\begin{itemize}
		\item Η φυσιολογική κίνηση (κλάση \texttt{0}) διασκορπίστηκε κυρίως στα Clusters 0, 1, 4, 5.
		\item Συγκεκριμένες επιθέσεις ομαδοποιήθηκαν μαζί με μεγάλη πυκνότητα. Για παράδειγμα, η κλάση \texttt{4} έχει ισχυρή παρουσία στο Cluster 1 ($791$ δείγματα) και στο Cluster 5 ($1739$ δείγματα).
	\end{itemize}
	
	\subsection{Τελικό Συμπέρασμα Σύγκρισης}
	Η σύγκριση των δύο προσεγγίσεων δείχνει ότι για το συγκεκριμένο πρόβλημα ανίχνευσης κυβερνοεπιθέσεων, η \textbf{Επιβλεπόμενη Μάθηση (και συγκεκριμένα το Random Forest)} είναι σαφώς ανώτερη και έτοιμη για πρακτική εφαρμογή, καθώς μπορεί να διαχειριστεί την ανισορροπία των κλάσεων μέσω κατάλληλων βαρών (\texttt{class\_weight='balanced'}). 
	
	Αντίθετα, η \textbf{Μη Επιβλεπόμενη Μάθηση} επηρεάζεται έντονα από την παρουσία ακραίων τιμών (outliers) και την κυριαρχία της κλάσης BENIGN. Ωστόσο, ο DBSCAN αποδείχθηκε πολύ πιο χρήσιμος από τον K-Means, καθώς η ικανότητά του να αναγνωρίζει αυθαίρετα σχήματα πυκνότητας του επέτρεψε να διακρίνει υποομάδες επιθέσεων (όπως της κλάσης \texttt{4}), καθιστώντας τον ένα καλό εργαλείο για την ανίχνευση νέων, άγνωστων μορφών ανωμαλιών στο δίκτυο.
\end{document}